\documentclass{article}
\title{Spring 2024 MATH 610 HW 4 Solutions}
\author{Jordan Hoffart}
\date{\today}
\usepackage{amsmath}
\usepackage{amsthm}
\usepackage{amsfonts}
\theoremstyle{plain}
\newtheorem{lemma}{Lemma}
\newtheorem{corollary}{Corollary}
\begin{document}
\maketitle
\section*{Problem 1}
Let $b > 0$ and $f \in C[0,1]$. Let $N > 0$ and $h = 1/(N+1)$ and $x_i = ih$.
Show that the downwind scheme
\begin{align*}
	-\frac{W_{i+1}-2W_i+W_{i-1}}{h^2} + b\frac{W_{i+1}-W_i}{h} & = f(x_i), \\
	W_0                                                        & = 0,      \\
	W_{N+1}                                                    & = 0,
\end{align*}
is conditionally stable, i.e. there is $\overline h > 0$ and $C > 0$ such that when $h < \overline h$,
\begin{lemma}\label{lemma1}
	Suppose that $1-bh > 0$.
	If
	\begin{align*}
		-\frac{W_{i+1}-2W_i+W_{i-1}}{h^2} + b\frac{W_{i+1}-W_i}{h} & = 0, \\
		W_0                                                        & = 0, \\
		W_{N+1}                                                    & = 0,
	\end{align*}
	then $W_i = 0$ for all $i$.
\end{lemma}
\begin{proof}
	Indeed, we have that
	\[\frac{1-bh}{2-bh}(W_i-W_{i+1}) + \frac{1}{2-bh}(W_i-W_{i-1})=0.\]
	Since the coefficients are positive and sum to $1$, this shows that $0$ is a convex combination of $W_i-W_{i+1}$ and $W_i-W_{i-1}$.
	Then either $W_i-W_{i+1} \leq 0 \leq W_i - W_{i-1}$ or $W_i-W_{i-1} \leq 0 \leq W_i - W_{i+1}$.
	In other words, $W_{i-1} \leq W_i \leq W_{i+1}$ or $W_{i+1} \leq W_i \leq W_{i-1}$.
	Furthermore, if $W_i > W_{i-1}$, then $W_{i+1} > W_i$, and if $W_i < W_{i-1}$, then $W_{i+1} < W_i$.
	This implies that all $W_i$ must be zero.
	Indeed, if $W_1 > 0 = W_0$, then $W_2 > W_1 > 0$, and we proceed inductively to conclude that $W_{N+1} > 0$, which is a contradiction.
	Similarly, if $W_1 < W_0 = 0$, then we proceed similarly to show $W_{N+1} < 0$, which is also a contradiction.
	Hence $W_1 = 0$.
	Then we repeat this argument again with $W_2$ through $W_N$ to conclude that all $W_i = 0$.
\end{proof}
\begin{lemma}\label{lemma2}
	Suppose that $1-bh > 0$.
	If
	\begin{align*}
		-\frac{W_{i+1}-2W_i+W_{i-1}}{h^2} + b\frac{W_{i+1}-W_i}{h} & = g(x_i) \leq 0, \\
		W_0                                                        & = 0,             \\
		W_{N+1}                                                    & = 0,
	\end{align*}
	then $W_i \leq 0$ for all $i$.
\end{lemma}
\begin{proof}
	Suppose instead that $W_i > 0$ for some $i$.
	Let $W_k = \max_j W_j$.
	Then necessarily $k \in \{1,\dots,N\}$, so we have that for the smallest such integer $k$,
	\[(1-bh)(W_k-W_{k+1}) + (W_k-W_{k-1}) = g(x_k).\]
	Both terms on the left side of the equality are non-negative while the right side is non-positive.
	Therefore, $W_k - W_{k-1} = 0$, so either $k$ is not the smallest integer in $\{1,\dots,N\}$ that attains the maximum or $k = 1$, in which case $W_k = 0$.
	Either conclusion results in a contradition, so $W_i \leq 0$ for all $i$.
\end{proof}
\begin{corollary}\label{corollary1}
	Suppose that $1-bh > 0$.
	If
	\begin{align*}
		-\frac{W_{i+1}-2W_i+W_{i-1}}{h^2} + b\frac{W_{i+1}-W_i}{h} & = g(x_i) \geq 0, \\
		W_0                                                        & = 0,             \\
		W_{N+1}                                                    & = 0,
	\end{align*}
	then $W_i \geq 0$ for all $i$.
\end{corollary}
\begin{proof}
	Let $W_i^- = -W_i$ and we have from the previous lemma applied to $W_i^-$ that $-W_i \leq 0$ for all $i$.
\end{proof}
\[\max_i|U_i| \leq C\max_i|f(x_i)|.\]
\begin{proof}
	Suppose that $1-bh > 0$.
	Then the square linear system from Lemma \ref{lemma1} admits only zero as its solution, which implies that for any function $g$, the system
	\begin{align*}
		-\frac{W_{i+1}-2W_i+W_{i-1}}{h^2} + b\frac{W_{i+1}-W_i}{h} & = g(x_i), \\
		W_0                                                        & = 0,      \\
		W_{N+1}                                                    & = 0,
	\end{align*}
	has a unique solution.
	In particular, there is a unique nontrivial solution $W^0$ to
	\begin{align*}
		-\frac{W_{i+1}^0-2W_i^0+W_{i-1}^0}{h^2} + b\frac{W_{i+1}^0-W_i^0}{h} & = 1, \\
		W_0^0                                                                & = 0, \\
		W_{N+1}^0                                                            & = 0.
	\end{align*}
	We now set $Z_i^+ = U_i - \max_{j=1,\dots,N}|f(x_j)|W_i^0$.
	Then $Z_i^+$ satisfies
	\begin{align*}
		-\frac{Z_{i+1}^+-2Z_i^++Z_{i-1}^+}{h^2} + b\frac{Z_{i+1}^+-Z_i^+}{h} & = f(x_i) - \max_{j=1,\dots,N}|f(x_j)| \leq 0, \\
		Z_0^+                                                                & = 0,                                          \\
		Z_{N+1}^+                                                            & = 0.
	\end{align*}
	We apply Lemma \ref{lemma2} to $g(x_i) = f(x_i) - \max_j|f(x_j)|$ and $W_i = Z_i^+$ to conclude that
	\[U_i \leq \max_j|f(x_j)|W_i^0.\]
	Now we let $Z_i^- = -U_i - \max_j|f(x_j)|W_i^0$ and we repeat the previous argument to conclude that
	\[-U_i \leq \max_j|f(x_j)|W_i^0.\]
	Combining these shows that
	\[\max_i|U_i|\leq (\max_iW_i^0)\max_j|f(x_j)|.\]
	The constant $C_h = \max_iW_i^0$ still currently depends on $h$.
	However, with a \emph{lot} of extra work, we can find a constant $C > 0$ such that $C_h < C$ for all $h$ sufficiently small (smaller than $1/b$ and possibly smaller than another complicated constant that depends on $b$).
	We will accept these details without proof, as the problem is complicated enough as is.
	This completes the proof.
\end{proof}

\section*{Problem 2}
Let $\widehat K$ be the unit reference triangle in $\mathbb R^2$ with vertices $v_1 = (0,0), v_2 = (1,0),$ and $v_3 = (0,1)$.
\subsection*{Part 1}
Show that there is a constant $C > 0$ such that
\[\|v\|_{L^2(\partial \widehat K)} \leq C(\|v\|_{L^2(\widehat K)} + \|\nabla v\|_{L^2(\widehat K)})\]
for all $v \in C^1(\widehat K)$.
\begin{proof}
	Let  $e_{ij}$ be the edge connecting $v_i$ to $v_j$.
	For $e_{12}$, we parameterize the triangle with the square as
	\[\varphi(s,t) = (ts,1-t)\]
	for $(s,t) \in [0,1]^2$.
	Then $e_{12}$ is just the image of $s \mapsto \varphi(s,1)$.
	Then the change of variables formulas for these maps are
	\[\int_K f = \int_{[0,1]^2} f\circ\varphi(s,t) t\,ds\,dt\]
	and
	\[\int_{e_{12}}f = \int_0^1 f(s,0)\,ds\]
	Then by the fundamental theorem of calculus and change-of-variables $(x,y) = (st, 1-t) = \varphi(s,t)$,
	\begin{align*}
		\|v\|_{L^2(e_{12})}^2 & = \int_0^1v(s,0)^2\,ds                                                                                   \\
		                      & = \int_0^1\int_0^1 \frac{d}{dt}\left(t^2v(\varphi(st,1-t))^2\right)\,dt\,ds                              \\
		                      & = \int_0^1\int_0^1 2tv(\varphi(st,1-t))^2 + 2tv(\varphi(s,t))\nabla v(\varphi(s,t))\cdot (ts,-t)\,dt\,ds \\
		                      & = 2\|v\|_{L^2(K)}^2 + 2(v, \nabla v \cdot (x,y-1))_{L^2(K)}                                              \\
		                      & \leq 2\|v\|_{L^2(K)}^2 + 2\|v\|_{L^2(K)}\|\nabla v \cdot (x,y-1)\|_{L^2(K)}                              \\
		                      & \leq 2\|v\|_{L^2(K)}^2 + 2\|v\|_{L^2(K)}\|\nabla v\|_{L^2(K)}                                            \\
		                      & \leq 2(\|v\|_{L^2(K)} + \|\nabla v\|_{L^2(K)})^2.
	\end{align*}
	By symmetry, we can show a similar bound for edge $e_{13}$ by considering the parameterization $\varphi(s,t) = (1-t,st)$.
	For edge $e_{23}$, we consider the parameterization $\varphi(s,t) = t(s,1-s)$.
	Then by following a similar argument to above, we have that
	\begin{align*}
		\frac{1}{\sqrt{2}}\|v\|_{L^2(e_{23})}^2 & = \int_0^1 v(s,1-s)^2\,ds                                \\
		                                        & = \int_0^1\int_0^1 \frac{d}{dt}(t^2v(ts,t(1-s)))\,dt\,ds \\
		                                        & \leq 2(\|v\|_{L^2(K)} + \|\nabla v\|_{L^2(K)})^2.
	\end{align*}
	Therefore,
	\begin{equation*}
		\|v\|_{L^2(\partial K)} \leq \sqrt{2(2 + \sqrt{2})}(\|v\|_{L^2(K)} + \|\nabla v\|_{L^2(K)}).
	\end{equation*}
	For convenience, we can replace the constant by $2\sqrt{2}$ or $4$.
\end{proof}

\subsection*{Part 2}
Let $K$ be any non-degenerate triangle in $\mathbb R^2$ and let $\widehat K$ be the reference triangle from above.
Find an affine transformation $F_K : \widehat K \to K$ as a function of the vertices of $K$.
Show that there are constants $C_1,C_2 > 0$ such that
\begin{equation*}
	C_1\|\xi\| \leq \|DF_K\xi\| \leq C_2\|\xi\|
\end{equation*}
for all $\xi \in \mathbb R^2$.
\begin{proof}
	Let $v_1,v_2,v_3$ be the vertices of $K$.
	Since $K$ is non-degenerate, $v_2 - v_1$ and $v_3 - v_1$ are linearly independent vectors in $\mathbb R^2$.
	We parametrize $K$ with $\widehat K$ by defining
	\begin{equation*}
		F_K(\widehat x, \widehat y) = v_1 + (v_2 - v_1)\widehat x + (v_3 - v_1)\widehat y.
	\end{equation*}
	Then
	\[DF_K = \begin{pmatrix} v_2 - v_1 & v_3 - v_1 \end{pmatrix},\]
	where $v_i - v_1$ are column vectors.
	Since the columns of $DF_K$ are linearly independent, $DF_K$ is an invertible matrix, so that the map $\xi \mapsto \|DF_K\xi\|$ is a norm on $\mathbb R^2$.
	Since all norms are equivalent on finite dimensional vector spaces, we are done.

\end{proof}

\subsection*{Part 3}
Let $h$ be the diameter of $K$.
Show that there is a constant $C > 0$ independent of $h$ such that for any $v \in C^1(K)$,
\[\|v\|_{L^2(\partial  K)} \leq C(h^{-1/2}\|v\|_{L^2( K)} + h^{1/2}\|\nabla v\|_{L^2( K)}).\]
\begin{proof}
	Let $E_i$ denote the edges of $K$ and $\widehat E_i$ denote the edges of $\widehat K$.
	Let $F_K$ be the affine transformation from the previous part that maps edge $\widehat E_i$ to edge $E_i$.
	Let $T_i$ be the restriction of $F_K$ to $\widehat E_i$.
	Then $T_i$ is also an affine transformation with $|\det DT_i| \leq h$ for each $i$.
	Then we have that
	\begin{equation*}
		\int_{E_i}f(s)^2\,ds \leq h\int_{\widehat E_i}f(T_i(\widehat s))^2\,d\widehat s
	\end{equation*}
	for any $f \in L^2(E_i)$.
	Then we have that
	\begin{align*}
		\|v\|_{L^2(\partial K)}^2 & = \sum_i \|v|_{E_i}\|_{L^2(E_i)}^2                      \\
		                          & \leq h\sum_i\|v|_{E_i}\circ T_i\|_{L^2(\widehat E_i)}^2 \\
		                          & = h\|v\circ F_K\|_{L^2(\partial\widehat K)}^2.
	\end{align*}
	Thus
	\begin{align*}
		\|v\|_{L^2(\partial K)} & \leq h^{1/2}\|v\circ F_K\|_{L^2(\partial\widehat K)}                                                   \\
		                        & \leq Ch^{1/2}\left(\|v\circ F_K\|_{L^2(\widehat K)} + \|\nabla(v\circ F_K)\|_{L^2(\widehat K)}\right).
	\end{align*}
	Since
	\begin{align*}
		\nabla(v\circ F_K) & = DF_K^T((\nabla v)\circ F_K)
	\end{align*}
	and
	\begin{equation*}
		\int_{\widehat K}f(\widehat x,\widehat y)\,d\widehat x\,d\widehat y = \frac{1}{|\det DF_K|}\int_Kf(x,y)\,dx\,dy
	\end{equation*}
	for any smooth $f$, we have that
	\begin{align*}
		\|\nabla (v\circ F_K)\|_{L^2(\widehat K)} \leq \frac{\|DF_K\|}{|\det DF_K|^{1/2}}\|\nabla v\|_{L^2(K)}
	\end{align*}
	where
	\[\|DF_K\| = \sup_{\xi \in \mathbb R^2\setminus\{0\}} \frac{\|DF_K\xi\|}{\|\xi\|}\]

	Since the previous proof does not give us more explicit constants that we need for this problem, let us show another proof that does this.
	Let $\rho_{\widehat K}$ be the diameter of the largest ball that fits inside of $\widehat K$.
	Then
	\begin{equation*}
		\|DF_K\| = \frac{1}{\rho_{\widehat K}}\sup_{\substack{\xi \in \mathbb R^2\\\|\xi\|=\rho_{\widehat K}}}\|DF_K\xi\|
	\end{equation*}
	For $\xi$ with $\|\xi\| = \rho_K$, we can write $\xi = (x_1,y_1) - (x_0,y_0)$ with $(x_i,y_i) \in \widehat K$.
	Then we have that $DF_K\xi = DF_K(x_1,y_1) = DF_K(x_0,y_0) = F_K(x_1,y_1) - F_K(x_0,y_0)$, so that $\|DF_K\xi\| \leq h_K$, the diameter of $K$.
	Thus $\|DF_K\| \leq \frac{h_K}{\rho_{\widehat K}}.$

	Furthermore, we have that
	\[|K| = \int_K1\,dx = \int_{\widehat K}|\det DF_K|\,d\widehat x = |\det DF_K||\widehat K| = |\det DF_K|/2\]
	so that $|\det DF_K| = 2|K|$.
	Since $K$ is a triangle with longest side length $h_K$, its height is $\alpha_Kh_K$ where $0 < \alpha_K < 1$.
	Therefore, $|K| = C_Kh_K^2$ for some constant $0 < C_K < 1$
	Thus,
	\begin{align*}
		\|\nabla (v\circ F_K)\|_{L^2(\widehat K)} \leq \frac{\|DF_K\|}{|\det DF_K|^{1/2}}\|\nabla v\|_{L^2(K)} \leq C_K\|\nabla v\|_{L^2(K)}.
	\end{align*}
	Similarly,
	\begin{align*}
		\|v\circ F_K\|_{L^2(K)} = \frac{1}{|\det DF_K|^{1/2}}\|v\|_{L^2(\widehat K)} \leq \frac{C_K}{h_K}\|v\|_{L^2(K)}.
	\end{align*}
	Combining our results gives us
	\begin{align*}
		\|v\|_{L^2(\partial K)} & \leq h^{1/2}\|v\circ F_K\|_{L^2(\partial\widehat K)}                                                  \\
		                        & \leq Ch^{1/2}\left(\|v\circ F_K\|_{L^2(\widehat K)} + \|\nabla(v\circ F_K)\|_{L^2(\widehat K)}\right) \\
		                        & \leq C_K\left(h^{-1/2}\|v\|_{L^2(K)} + h^{1/2}\|\nabla v\|_{L^2(K)}\right)
	\end{align*}
	where the final constant $C$ does not depend on the diameter of $K$ and does not depend on $v$ but may depend on the height of $K$ as well as the reference element $\widehat K$.

\end{proof}




\end{document}
